\section{Experiments}
\label{sec:experiments}

This section evaluates Toollery as a self-verified query retrieval framework for
skill and atomic-skill selection. We use
\emph{skill} as the general interface unit and \emph{atomic skill} for the original
tool-style APIs whose behavior can be described by a name, natural-language
description, schema, and optional manual. The central question is not whether a
larger prompt can contain more generated text, but whether self-verified usage
queries can build a better retrieval index while keeping the final LLM decision
prompt compact. Unless otherwise noted, generated and verified queries are used
only inside the retrieval index. The final LLM selector receives only the top-$k$
compact candidate skill specifications, not the generated query pool.

\subsection{Research Questions}
\label{subsec:rq}

\paragraph{RQ1: Selection accuracy.}
Does Toollery improve skill selection over raw full-text retrieval, dense
retrieval, reranking, SkillRouter-style pipelines, and RAG-style baselines?

\paragraph{RQ2: Benefit of query augmentation.}
How much of the gain comes from adding generated usage queries to the index, and
how much comes specifically from Toollery's round-trip self-verification?

\paragraph{RQ3: Skill and atomic-skill generality.}
Does the same retrieval and verification design work on both SkillBench-style
skill routing tasks and the original atomic tool benchmarks such as HammerBench,
xLAM, BFCL, and ScaleTool?

\paragraph{RQ4: Efficiency under large candidate pools.}
Can Toollery reduce final-prompt tokens, LLM selection latency, and monetary cost
by retrieving a compact subset before the LLM sees the candidates, while
preserving or improving selection quality?

\paragraph{RQ5: Robustness to distractors and scale.}
How does performance change as the candidate skill pool grows from small
benchmark pools to large mixed pools with many semantically similar distractors?

\subsection{Benchmarks}
\label{subsec:benchmarks}

\paragraph{SkillBench / SkillRouter-compatible track.}
This track evaluates Toollery on skill routing datasets where each instance
contains a user request and one or more gold skills. To make the comparison
compatible with SkillRouter-style evaluation, we report retrieval-oriented
metrics such as Hit@$k$, Recall@$k$, MRR, and nDCG whenever the annotations
support ranked relevance. When executable task environments or validated final
answers are available, we additionally report final task success after the LLM
selects from the retrieved candidate subset.

\paragraph{Atomic tool/skill track.}
This track keeps the original Toollery experiments on HammerBench, xLAM, BFCL,
and ScaleTool, but reframes each tool as an atomic skill. Each atomic skill is
represented by its original metadata, including name, description, schema,
manual, and example fields when available. We retain the original evaluation
metrics from the paper for comparability, and add retrieval metrics plus
efficiency measurements to isolate the effect of candidate pruning.

\paragraph{Large-pool and distractor settings.}
For both tracks, we create controlled scale settings by increasing the number of
candidate skills available at inference time. Pool sizes should include the
native benchmark pool and larger mixed pools, for example $N \in \{100, 500,
1{,}000, 5{,}000, 10{,}000\}$ when data availability allows. We also create
hard-distractor pools by mixing skills with overlapping names, descriptions,
schemas, or task domains. This setting tests whether verified queries improve
grounding beyond surface-level name and description matching.

\subsection{Baselines and Settings}
\label{subsec:baselines}

We separate method families from information settings. This avoids treating raw
metadata retrieval as if it were a Toollery variant. Toollery is defined by query-level
retrieval plus self-verification; raw metadata retrieval is a baseline setting,
not a Toollery variant.

\paragraph{Information settings.}
\textbf{Raw Baselines} use only the original skill metadata, such as name,
description, schema, manual, body, and examples. \textbf{Query-Augmented
Baselines} add the same generated usage queries to each baseline's searchable
index to test whether query generation alone improves retrieval. \textbf{
Unverified Toollery} retrieves over generated usage queries and aggregates
query-level matches to skill-level candidates, but does not apply round-trip
self-verification. \textbf{Full Toollery} generates candidate usage queries,
filters them through round-trip self-verification, indexes only verified queries,
retrieves matching verified queries at inference time, aggregates them into a
top-$k$ skill subset, and sends only compact candidate skill specifications to
the final LLM selector.

\paragraph{Retrieval baselines.}
We include BM25 over raw skill text, dense retrieval over raw skill text, and a
retrieve-and-rerank pipeline that mirrors the SkillRouter family when model
access and training data are available. For the SkillBench-compatible track, the
SkillRouter comparison should use the official encoder and reranker setup where
possible. If reproducing the full training recipe is infeasible, we report a
frozen off-the-shelf dense retriever plus reranker and label it as an adapted
SkillRouter-style baseline rather than the official SkillRouter model.

\paragraph{RAGAnything baseline.}
RAGAnything is adapted as a generic retrieval-over-documentation baseline. For
raw settings, it indexes each skill's original documentation fields. For
query-augmented settings, it receives the same generated queries used by the
other baselines. This comparison tests whether Toollery's gain comes from
generic document retrieval or from the verified query-to-skill retrieval
structure.

\paragraph{LLM selector baselines.}
For small pools, we include a direct LLM selector that sees all candidate skill
specifications. For larger pools, the all-candidate prompt may exceed context or
be economically unrealistic; in that case we report the maximum feasible pool
size and compare against retrieval-preselected LLM selection. The final selector
model, prompt template, temperature, maximum output length, and candidate format
must be fixed across methods.

\subsection{Metrics}
\label{subsec:metrics}

\paragraph{Retrieval and routing metrics.}
For single-gold tasks, we report Hit@$k$ for $k \in \{1, 3, 5, 10, 20\}$ and
MRR. For multi-gold tasks, we report Recall@$k$ and nDCG@$k$ when graded or
multi-relevance labels are available. We also report candidate-set recall before
the final LLM selection step, because Toollery's first-stage objective is to
retain the gold skill inside a compact subset.

\paragraph{Final task metrics.}
On executable or answer-validated benchmarks, we report final task success,
exact match, pass rate, or the benchmark's official score. The key distinction is
between retrieval success and end-to-end success: a method can retrieve the gold
skill but fail if the final LLM chooses incorrectly or formats the call
incorrectly.

\paragraph{Efficiency metrics.}
We measure the number of candidates sent to the final LLM, final prompt tokens,
final completion tokens, retrieval latency, reranking latency, final LLM
selection latency, total inference latency, and estimated cost per 1,000 user
queries. We also report compression ratios relative to the direct all-candidate
LLM selector:
\[
\mathrm{TokenCompression}
= \frac{\mathrm{PromptTokens}_{\mathrm{all}}}
{\mathrm{PromptTokens}_{\mathrm{method}}},
\quad
\mathrm{LatencyCompression}
= \frac{\mathrm{LLMLatency}_{\mathrm{all}}}
{\mathrm{LLMLatency}_{\mathrm{method}}}.
\]
When all-candidate prompting is infeasible for large pools, we use the largest
feasible direct-LLM pool as an upper-bound reference and clearly mark larger
pool comparisons as retrieval-only or retrieval-plus-LLM-subset evaluation.

\paragraph{Verification metrics.}
For the offline query construction stage, we report the number of generated
queries per skill, the self-verification acceptance rate, the average number of
verified queries per skill, the percentage of skills with zero verified queries,
offline construction cost, and offline construction latency. These metrics
separate one-time indexing cost from per-query serving cost.

\subsection{Ablations}
\label{subsec:ablations}

\paragraph{Generated queries without verification.}
We compare Raw Baselines, Query-Augmented Baselines, Unverified Toollery, and
Full Toollery. This is the main ablation requested by the current paper
positioning. If generated queries help every method, Query-Augmented Baselines
should improve over Raw Baselines. If round-trip self-verification removes noisy
or misleading generated queries, Full Toollery should improve over Unverified
Toollery, especially in hard-distractor settings.

\paragraph{Verification strictness.}
We vary the round-trip self-verification threshold or acceptance rule. Candidate
queries can be accepted if the verifier recovers the source skill at top-1,
top-3, or top-5, or if the verifier assigns a score above a calibrated
threshold. This ablation traces the tradeoff between coverage and precision in
the verified query index.

\paragraph{Number of generated queries.}
We vary the number of generated candidate queries per skill before verification,
for example $m \in \{5, 10, 20, 40\}$. The expected analysis is whether larger
query pools improve recall until verification or retrieval becomes saturated.

\paragraph{Candidate subset size.}
We vary the number of compact candidates sent to the final LLM, for example
$k \in \{1, 3, 5, 10, 20\}$. This ablation directly measures the quality-cost
frontier: smaller $k$ should reduce tokens and LLM latency, while larger $k$
should increase candidate-set recall and potentially final task success.

\paragraph{Index field ablation.}
We compare indexes over name-only, name plus description, full raw text, raw text
plus generated queries, generated queries only, and verified queries only. This
helps position Toollery against SkillRouter-style full-text retrieval and tests
whether verified query text captures usage semantics that raw documentation
misses.

\paragraph{Aggregation ablation.}
Because Toollery retrieves at the query level and then aggregates to the skill
level, we compare aggregation rules such as max score, mean top-$r$ score, sum of
top-$r$ scores, reciprocal-rank fusion, and learned lightweight aggregation if
training data are available.

\subsection{Efficiency, Token, and Latency Analysis}
\label{subsec:efficiency}

The efficiency claim should be evaluated as a first-class result, not a side
note. Toollery may build a larger offline retrieval index, but that index is not
included in the final LLM prompt. At serving time, the system retrieves verified
queries, maps them back to a small skill subset, and sends only compact
candidate specifications to the final LLM selector. Therefore the online token
cost scales with $k$, not with the full number of indexed queries or the full
skill pool.

For each method and pool size, we log:
\begin{itemize}
    \item \textbf{Index size:} total indexed documents, average indexed entries
    per skill, and average indexed tokens per skill.
    \item \textbf{Final prompt size:} number of candidates sent to the LLM and
    prompt tokens after candidate pruning.
    \item \textbf{Latency:} retrieval latency, reranking latency, final LLM
    selection latency, and total latency.
    \item \textbf{Cost:} estimated per-query LLM cost and total cost per 1,000
    queries, separating offline indexing cost from online serving cost.
    \item \textbf{Quality-efficiency frontier:} final success or Hit@$k$ versus
    prompt tokens and total latency.
\end{itemize}

The main efficiency comparison should include three practically important
systems: direct all-candidate LLM selection, SkillRouter/RAG-style retrieval plus
LLM selection over a subset, and Full Toollery retrieval plus LLM selection over
a subset. This distinguishes prompt compression from retrieval quality.

\subsection{Expected Tables}
\label{subsec:expected-tables}

\begin{table*}[t]
\centering
\small
\begin{tabular}{lcccccc}
\hline
\textbf{Setting} &
\textbf{Raw fields} &
\textbf{Generated queries} &
\textbf{Self-verification} &
\textbf{Retrieval unit} &
\textbf{Final LLM input} &
\textbf{Purpose} \\
\hline
Raw Baselines &
\checkmark & -- & -- &
Skill document &
Top-$k$ compact specs &
Raw metadata retrieval \\
Query-Augmented Baselines &
\checkmark & \checkmark & -- &
Skill document or query-augmented document &
Top-$k$ compact specs &
Effect of generation alone \\
Unverified Toollery &
Optional & \checkmark & -- &
Generated query &
Top-$k$ compact specs &
Query-level retrieval without filtering \\
Full Toollery &
Optional & \checkmark & \checkmark &
Verified query &
Top-$k$ compact specs &
Verified query index plus compact LLM selection \\
\hline
\end{tabular}
\caption{Experimental settings. Generated or verified queries are indexed for
retrieval but are not copied into the final LLM prompt. The final selector sees
only compact specifications for the retrieved top-$k$ skills.}
\label{tab:settings-overview}
\end{table*}

\begin{table*}[t]
\centering
\small
\begin{tabular}{llcccccc}
\hline
\textbf{Track} &
\textbf{Method} &
\textbf{Hit@1} &
\textbf{Hit@5} &
\textbf{Recall@10} &
\textbf{MRR} &
\textbf{nDCG@10} &
\textbf{Final Success} \\
\hline
SkillBench & BM25 (Raw) & -- & -- & -- & -- & -- & -- \\
SkillBench & Dense Retriever (Raw) & -- & -- & -- & -- & -- & -- \\
SkillBench & SkillRouter / SkillRouter-style & -- & -- & -- & -- & -- & -- \\
SkillBench & RAGAnything (Raw) & -- & -- & -- & -- & -- & -- \\
SkillBench & RAGAnything (+ generated queries) & -- & -- & -- & -- & -- & -- \\
SkillBench & Unverified Toollery & -- & -- & -- & -- & -- & -- \\
SkillBench & Full Toollery & -- & -- & -- & -- & -- & -- \\
\hline
Atomic skills & BM25 (Raw) & -- & -- & -- & -- & -- & -- \\
Atomic skills & Dense Retriever (Raw) & -- & -- & -- & -- & -- & -- \\
Atomic skills & RAGAnything (Raw) & -- & -- & -- & -- & -- & -- \\
Atomic skills & RAGAnything (+ generated queries) & -- & -- & -- & -- & -- & -- \\
Atomic skills & Unverified Toollery & -- & -- & -- & -- & -- & -- \\
Atomic skills & Full Toollery & -- & -- & -- & -- & -- & -- \\
\hline
\end{tabular}
\caption{Main selection and end-to-end results. Use SkillBench-compatible
metrics on the skill-routing track and the original benchmark official score for
the atomic-skill track when available. All entries are placeholders until the
experiments are run.}
\label{tab:main-selection-results}
\end{table*}

\begin{table*}[t]
\centering
\small
\begin{tabular}{lccccccc}
\hline
\textbf{Method / Ablation} &
\textbf{Gen. queries / skill} &
\textbf{Accepted / skill} &
\textbf{Accept rate} &
\textbf{Zero-query skills} &
\textbf{Hit@5} &
\textbf{Final Success} &
\textbf{Offline Cost} \\
\hline
Raw Baseline & 0 & 0 & -- & -- & -- & -- & -- \\
Query-Augmented Baseline & $m$ & $m$ & 100\% & 0\% & -- & -- & -- \\
Unverified Toollery & $m$ & $m$ & 100\% & 0\% & -- & -- & -- \\
Full Toollery (top-1 verification) & $m$ & -- & -- & -- & -- & -- & -- \\
Full Toollery (top-3 verification) & $m$ & -- & -- & -- & -- & -- & -- \\
Full Toollery (score threshold) & $m$ & -- & -- & -- & -- & -- & -- \\
\hline
\end{tabular}
\caption{Verification ablation. This table isolates the contribution of
round-trip self-verification from query generation itself. Offline cost should be
reported separately from online inference cost.}
\label{tab:verification-ablation}
\end{table*}

\begin{table*}[t]
\centering
\small
\begin{tabular}{lcccccccc}
\hline
\textbf{Method} &
\textbf{Pool size} &
\textbf{Cand. to LLM} &
\textbf{Prompt toks} &
\textbf{Tok. comp.} &
\textbf{Retr. ms} &
\textbf{Rerank ms} &
\textbf{LLM ms} &
\textbf{Cost / 1k} \\
\hline
Direct all-candidate LLM & -- & -- & -- & 1.0$\times$ & -- & -- & -- & -- \\
BM25 + LLM subset & -- & $k$ & -- & -- & -- & -- & -- & -- \\
SkillRouter-style + LLM subset & -- & $k$ & -- & -- & -- & -- & -- & -- \\
RAGAnything + LLM subset & -- & $k$ & -- & -- & -- & -- & -- & -- \\
Unverified Toollery + LLM subset & -- & $k$ & -- & -- & -- & -- & -- & -- \\
Full Toollery + LLM subset & -- & $k$ & -- & -- & -- & -- & -- & -- \\
\hline
\end{tabular}
\caption{Token, latency, and cost compression. Generated and verified queries
increase the retrieval index size but do not enter the final LLM prompt. The
online LLM prompt contains only the compact specifications of the retrieved
top-$k$ candidates.}
\label{tab:token-latency-compression}
\end{table*}

\begin{table*}[t]
\centering
\small
\begin{tabular}{lcccccc}
\hline
\textbf{Method} &
\textbf{$N=100$} &
\textbf{$N=500$} &
\textbf{$N=1{,}000$} &
\textbf{$N=5{,}000$} &
\textbf{$N=10{,}000$} &
\textbf{Avg. prompt toks} \\
\hline
BM25 (Raw) & -- & -- & -- & -- & -- & -- \\
Dense Retriever (Raw) & -- & -- & -- & -- & -- & -- \\
RAGAnything (+ generated queries) & -- & -- & -- & -- & -- & -- \\
SkillRouter-style & -- & -- & -- & -- & -- & -- \\
Unverified Toollery & -- & -- & -- & -- & -- & -- \\
Full Toollery & -- & -- & -- & -- & -- & -- \\
\hline
\end{tabular}
\caption{Scaling curve under increasing candidate-pool size. Each cell can
report Hit@5, Recall@10, or final success depending on the benchmark. A companion
figure should plot quality versus pool size and quality versus final prompt
tokens.}
\label{tab:scaling-curve}
\end{table*}

\subsection{Experiments to Run}
\label{subsec:experiments-to-run}

\begin{enumerate}
    \item Build the SkillBench / SkillRouter-compatible evaluation pipeline:
    normalize skill metadata, define gold labels, implement Hit@$k$,
    Recall@$k$, MRR, nDCG, and final task success when supported.
    \item Re-run the original Toollery benchmark suite on HammerBench, xLAM,
    BFCL, and ScaleTool with terminology changed from tool to atomic skill.
    \item Implement Raw Baselines for BM25, dense retrieval, reranker retrieval,
    RAGAnything, and direct LLM selection where context length permits.
    \item Generate the shared query pool for every skill. Store the generated
    queries once and reuse the identical pool for Query-Augmented Baselines,
    Unverified Toollery, and Full Toollery.
    \item Run Query-Augmented Baselines by adding the shared generated queries to
    each baseline's searchable index without using Toollery's round-trip
    verification filter.
    \item Run Unverified Toollery with query-level retrieval and skill-level
    aggregation, using generated queries but no self-verification.
    \item Run Full Toollery with round-trip self-verification, a verified query
    index, query-level retrieval, skill-level aggregation, and compact top-$k$
    final LLM selection.
    \item Add RAGAnything on the original atomic-skill benchmarks in both raw and
    query-augmented forms.
    \item If feasible, evaluate official SkillRouter on SkillBench. If not
    feasible, clearly label the reproduced comparison as SkillRouter-style and
    document the encoder, reranker, and training or zero-shot setup.
    \item Measure token count, retrieval latency, reranking latency, final LLM
    latency, total latency, and estimated cost for each method and pool size.
    \item Run ablations over verification strictness, number of generated
    queries per skill, candidate subset size $k$, index fields, and aggregation
    rule.
    \item Run scale and hard-distractor evaluations by increasing candidate pool
    size and adding semantically similar skills.
\end{enumerate}

\subsection{Implementation Notes}
\label{subsec:implementation-notes}

\paragraph{Fair query augmentation.}
The generated query pool must be identical across Query-Augmented Baselines,
Unverified Toollery, and Full Toollery before verification. This controls for the
benefit of generation itself. Full Toollery may remove queries through
self-verification, but it should not receive additional private generated
queries that are unavailable to the ablations.

\paragraph{No generated queries in final prompts.}
For all Toollery variants and query-augmented baselines, final LLM prompts should
contain only compact candidate skill specifications. If a baseline uses
generated queries in the final prompt, report it as a separate prompt-expansion
baseline because it changes the token budget and weakens the compression claim.

\paragraph{Prompt and model control.}
Use the same final selector model, system prompt, candidate formatting,
temperature, and decoding parameters for all methods. The only allowed
difference should be the retrieved candidate subset and its order.

\paragraph{Latency measurement.}
Measure retrieval, reranking, and final LLM latency separately. Report mean,
median, and p95 if possible. Offline generation and verification costs should be
reported separately from online serving latency.

\paragraph{Index construction.}
Store raw skill documents, generated queries, and verified queries as separate
index fields or separate indexes. This makes it possible to run field ablations
without regenerating data.

\paragraph{RAGAnything adaptation.}
Document exactly how each skill is serialized into RAGAnything documents and how
retrieved chunks are mapped back to skill IDs. If RAGAnything retrieves chunks
rather than skills, aggregate chunk scores to skill scores before top-$k$ final
LLM selection.

\paragraph{Statistical reporting.}
Run at least three seeds for query generation and LLM selection when budget
allows. Report confidence intervals or bootstrap standard errors for main
results. If generation is expensive, fix one generated query pool and vary only
retrieval and final-selection seeds, while explicitly stating this limitation.

\paragraph{Failure analysis.}
Sample failures where Full Toollery loses to Query-Augmented Baselines or raw
retrieval. Categorize them into missing generated coverage, over-strict
verification, ambiguous gold labels, schema-level mismatch, and final LLM
selection error. This analysis will help explain whether verification improves
precision at the cost of recall.
